%##########################################################################
% CHAPTER 5: DISCUSSION
%##########################################################################
\chapter{Discussion}

\section{Chapter Overview}

This chapter steps back from the numbers and themes presented in Chapter~4 and asks what they mean. It interprets the quantitative and qualitative findings in relation to the research questions and the existing literature, identifies the original contributions of the work, and, with equal candour, acknowledges where the system falls short and why.

\section{Interpretation of Quantitative Findings}

The central prediction of this research, that PPO-driven narrative adaptation would produce measurably higher Flow experience, rests on a theoretical chain with several links. The player modelling module must accurately infer the player's psychological state from noisy behavioural telemetry. The PPO agent must select a narratively appropriate action in response to that inferred state. The LLM must generate content that faithfully executes the selected action while remaining grounded in the game's lore. And the player must actually experience the resulting narrative shift as meaningful: not as noise, not as interruption, but as a story that responds to them. If any link in this chain fails, the effect on the dependent variable (GEQ Flow subscale) will be attenuated or absent.

The GEQ Flow results ($d$ = 0.65, $p$ = .027) provide partial support for H1. The medium effect size indicates that PPO-driven narrative adaptation produced a meaningful improvement in self-reported Flow, though the result narrowly missed the Bonferroni-corrected threshold ($\alpha_{\text{adj}}$ = .0167). This outcome is interpretable in two complementary ways. On one hand, it demonstrates that narrative tone selection can influence the psychological state that Csikszentmihalyi's (\citeyear{csikszentmihalyi1990flow}) theory places at the centre of optimal experience, extending the scope of Flow-targeted adaptation beyond the mechanical DDA interventions reviewed in Section~2.4. On the other hand, the marginal significance under correction suggests that the effect, while real, is not yet robust enough to survive the conservative multiple-comparison standard that the study design imposed.

The convergent evidence from the behavioural metrics substantially strengthens this interpretation, because behavioural measures are immune to the demand characteristics that threaten self-report scales. The dialogue read ratio was substantially higher in the Experimental condition ($M$ = 0.72 vs.\ 0.58, $d$ = 0.87), indicating that participants were not merely \textit{reporting} higher Flow retrospectively but were also \textit{demonstrably} more engaged with the narrative content during play. Crucially, participants did not know that their dialogue reading behaviour was being measured, so this difference cannot be attributed to response bias. Taken together, these behavioural indicators provide demand-characteristic-free corroboration of the self-report findings: the adaptive narrative did not merely change what participants said about their experience; it changed what they actually did during play.

The GEQ Immersion result ($d$ = 0.38, $p$ = .187) did not reach significance. The small effect size suggests that a 30-minute session with a prototype-quality testbed game may be insufficient for narrative adaptation to influence the broader, multisensory construct of immersion. Immersion, as measured by the GEQ, encompasses factors such as visual richness, spatial presence, and loss of awareness of the real world, many of which depend on production values (art, audio, level design) that the testbed does not possess. Narrative tone alone may not be a strong enough lever to move this particular subscale in a low-fidelity context.

The NPS-3 result ($d$ = 0.84, $p$ = .005) is the strongest finding. The large effect size and survival under Bonferroni correction indicate that participants in the Experimental condition clearly perceived the narrative as more responsive to their behaviour. This is consistent with the system's design: the PPO agent selects among eight distinct narrative actions based on the player's state, and the resulting variation in tone and content is, at least for the majority of participants, perceptible. The NPS-3 result also aligns with the qualitative Theme~1 (``The story pays attention''), in which 8 of 10 interviewed Experimental participants described a sense that the game's narrative was responding to them.

The pattern across the three hypotheses, strong personalization, moderate Flow, weak immersion, is internally coherent. The system directly manipulates narrative content in response to player state (captured by NPS-3), and this manipulation has a downstream effect on Flow (the PPO optimisation target), but does not reach far enough to influence the broader construct of immersion, which depends on factors outside the system's control.

\section{Themes from Interviews}

The four themes identified through Reflexive Thematic Analysis illuminate the mechanisms behind the quantitative results and surface nuances that the structured scales could not capture.

Theme~1 (``The story pays attention''), reported by 8 of 10 Experimental interviewees, provides direct qualitative support for the NPS-3 finding. Participants did not describe personalization in terms of explicit branching choices; instead, they described a more ambient sense that the narrative world was responsive to how they were playing. The metaphor of being ``noticed'' or ``seen'' by the game recurred across several interviews. This is precisely the kind of experience the system was designed to produce, and its prevalence among Experimental participants suggests that the adaptation pipeline's effect is perceptible at the level of subjective experience, not merely at the level of aggregated scale scores.

Theme~2 (``Tone matching'') connects directly to the PPO agent's action selection. Six Experimental participants described specific moments where the narrative tone aligned with their gameplay intensity: supportive, urgent dialogue during difficult encounters; lore-rich, atmospheric content during exploratory phases. This alignment is the behavioural signature of effective PPO action selection, and the fact that players noticed and valued it suggests that the reward function's Flow-targeting objective translates into perceptible narrative quality.

Theme~3 (``Generic narrative as wallpaper''), prevalent among Control participants, illuminates the mechanism behind the dialogue read ratio difference. The static narrative was not perceived as poorly written; it was perceived as irrelevant. Seven of 10 Control interviewees described a gradual disengagement from narrative content over the session, consistent with the 58\% dialogue read ratio. This theme suggests that the absence of personalization does not produce active dissatisfaction but rather a quiet disengagement, a finding that aligns with Short's (\citeyear{short2019procedural}) critique of procedural generation without player intent modelling.

Theme~4 (``Occasional dissonance'') identifies a concrete failure mode. Four Experimental participants noted moments where the narrative tone felt mismatched, most commonly receiving ``mystery'' or ``lore reward'' content during or immediately after intense combat. This pattern suggests that the PPO agent's action selection occasionally lags behind rapid state transitions. The 5-second telemetry cycle means that a player who transitions from ANXIETY to FLOW within a few seconds may receive a narrative action selected for the prior state. This latency-induced mismatch is a design limitation rather than a policy failure, and it suggests that future work should explore event-triggered telemetry updates in addition to the fixed 5-second cycle.

Notably, no participants in either condition described the narrative as obviously machine-generated or ``AI-sounding.'' This absence is consistent with the finding of \citet{ligs2025chi} that players cannot reliably distinguish LLM-generated narrative from human-authored content when the model is appropriately constrained. The RAG grounding and structured prompt appear to have been sufficient to maintain the fictional frame, even when the generated content occasionally misfired in tone.

\section{Novelty and Contributions}

The original contributions of this work are six in number, and they range from the practical (a reusable mapping table) to the conceptual (reframing what reinforcement learning in games is for).

\textbf{Contribution 1: Telemetry-derived behavioural motivation proxies aligned to the Hexad framework.} Every prior system that has used Hexad player types for game adaptation has relied on a pre-play questionnaire \citep{tondello2016hexad}, a static snapshot that cannot update during a session. This system computes six-dimensional continuous scores directly from behavioural telemetry, enabling within-session updating without any self-report instrument. The resulting profiles are best understood as \textit{behavioural motivation proxies}: Hexad-analogous scores that capture observable play patterns theoretically associated with each Hexad dimension, rather than validated measures of the underlying psychological constructs. The construct validity of these proxies against questionnaire-derived ground truth has not been established (see Section~5.5), and future work should administer the standard Hexad questionnaire alongside the telemetry-derived profiles to quantify the correlation. The telemetry-to-Hexad mapping (Table~\ref{tab:hexad}) is published in full and is reusable by other researchers and developers working with different games and different telemetry schemas.

\textbf{Contribution 2: RAG-grounded dynamic narrative.} PANGeA \citep{todd2024pangea} grounds LLM narrative in prompt-delivered context, but that context is assembled at design time and remains fixed across players and sessions. In the present system, the lore passages delivered to the LLM are selected at runtime based on the player's current location, activity, and the narrative action chosen by the PPO agent. The grounding is therefore dynamic rather than static: the same player in the same location at different Flow states will receive different lore context, producing different narrative output.

\textbf{Contribution 3: Episodic session memory in the LLM prompt.} The importance-weighted episodic memory, inspired by EM-LLM \citep{fountas2024episodic}, gives the LLM access to a compressed narrative history so that generated content can reference prior events, maintain character continuity, and avoid contradicting what the player has already experienced. To the author's knowledge, this is the first application of episodic memory compression to real-time game narrative generation.

\textbf{Contribution 4: PPO for narrative action selection (feasibility demonstration).} Reinforcement learning has been applied extensively to mechanical game adaptation (adjusting enemy health, spawn rates, and resource drops) but not, to the author's knowledge, to selecting \textit{narrative actions}. Framing narrative tone selection as an RL problem is a conceptual step beyond both rule-based narrative adaptation and mechanical DDA. The present implementation demonstrates the \textit{feasibility} of this framing: the PPO agent was trained on a simulated MDP with hand-specified transition dynamics (Section~3.6.1) and deployed without online updating, so the learned policy reflects the simulator's assumptions about player behaviour rather than empirically observed transitions. The contribution is therefore the framing and the proof-of-concept integration, not a claim that the specific learned policy is optimal. Online RL with real player data (Section~6.3) would be needed to validate the policy itself.

\textbf{Contribution 5: Flow-based MDP reward.} The reward function directly operationalises Csikszentmihalyi's (\citeyear{csikszentmihalyi1990flow}) challenge--skill balance theory, giving the adaptation objective a principled theoretical grounding rather than an ad hoc engagement metric. The graduated penalty structure (APATHY $>$ ANXIETY $>$ BOREDOM) reflects the severity ordering of disengagement states, and the shaping rewards (streak bonus, transition bonus, repetition penalty) encode design-level knowledge about what good adaptation looks like.

\textbf{Contribution 6: Complete closed-loop integration.} The five modules (telemetry capture, feature extraction and player typing, RL action selection, RAG-grounded LLM generation, and in-game content injection) run end-to-end within a single deployable Unity plugin. Individual pieces of this pipeline have appeared in prior work; what has not appeared is the closed loop. The player's behaviour influences the narrative, and the narrative influences the player's behaviour, creating a feedback cycle that is the defining characteristic of a genuinely adaptive system.

\section{Limitations}

No system is stronger than its weakest assumption, and this one has several that warrant explicit acknowledgement.

\textbf{Simulation fidelity.} The PPO agent is trained on a simulated environment with hand-specified transition probabilities: a 0.75 probability that a ``helpful'' action will move the player to FLOW, and a 0.25 probability that it will not. Real player behaviour is substantially more complex, more stochastic, and more history-dependent. The simulation's Markovian assumption, that the next state depends only on the current state and action, not on the sequence of states that preceded it, is almost certainly violated in practice. A player who has been in ANXIETY for five minutes is not in the same psychological position as one who just entered ANXIETY from FLOW, even if the telemetry-derived state vector is identical. The trained policy may therefore underperform in the real game environment relative to its performance in simulation.

\textbf{Rule-based Hexad profiling.} The mapping from telemetry signals to Hexad dimensions is based on theoretical alignment with Tondello et al.'s (\citeyear{tondello2016hexad}) motivation descriptions, not on empirical validation against questionnaire-derived ground truth. This means that the system's internal Hexad profiles may diverge from what a validated questionnaire would produce for the same player. Systematic misclassification is possible: a player who moves slowly due to hardware input lag or unfamiliarity with the controls, rather than deliberate exploration, will register a high \texttt{explore\_rate} and be profiled as an Explorer. The system has no way to distinguish intentional from incidental behaviour.

\textbf{LLM generation quality.} Phi-3.5 Mini is a 3.8-billion-parameter model, capable for its size but fundamentally constrained relative to larger frontier models. The quality of generated narrative (its fluency, creativity, emotional range, and consistency of character voice) is bounded by the model's capacity. Larger models would likely produce more compelling output, but they would also introduce latency, memory, and deployment constraints that conflict with the real-time requirement.

\textbf{Single-session study.} The evaluation captures a single 30-minute session per participant. Any effects of narrative personalization that require multiple sessions to develop (deepening attachment to characters, evolving expectations, long-term engagement) are invisible to this design. A longitudinal study would be needed to assess cumulative impact.

\textbf{Testbed ecological validity.} The testbed game is a simplified prototype built on a generic RPG framework. It lacks the production polish, narrative complexity, and mechanical depth of a commercially developed title. Findings obtained in this controlled environment may not generalise to games with more complex narrative architectures, larger worlds, or longer play sessions.

\textbf{Sample size.} With $n = 25$ per condition, the study is powered to detect medium-to-large effects ($d \geq 0.5$) but would miss small effects entirely. The convenience sample, drawn from university students and gaming society members, also limits demographic diversity and may not represent the broader population of RPG players.

\section{Threats to Validity}

\textbf{Internal validity.} The primary threat is demand characteristics. Participants who infer that they are in the experimental condition, perhaps because they notice that the narrative seems to be responding to their behaviour, may respond more favourably on the post-session questionnaires, inflating the apparent effect. The blind design (participants are not told which condition they are in) mitigates this threat, but does not eliminate it entirely. A more robust mitigation would be a double-blind design in which the researcher administering the session is also unaware of condition assignment, but the single-researcher study setup makes this impractical.

\textbf{Construct validity.} The GEQ Flow subscale measures retrospective self-report of Flow: what the participant \textit{remembers} feeling during the session, not what they \textit{actually} felt moment to moment. Peak-end effects and recency bias may distort the score. The absence of a physiological measure (e.g., galvanic skin response, heart rate variability) limits confidence that the GEQ score accurately reflects the within-session psychological state. The custom NPS-3 scale presents a more fundamental construct validity concern: although its internal consistency is acceptable ($\alpha = .76$), it has not undergone independent validation, factor analysis, or convergent/discriminant validity testing against established personalization or agency measures. The NPS-3 result ($d = 0.84$) is therefore best interpreted as evidence of perceived narrative responsiveness rather than as a validated measurement of ``narrative personalization'' as a psychological construct. Future work should either validate the NPS-3 against existing instruments or replace it with an established measure.

\textbf{External validity.} The testbed game is a simplified prototype, and the participants are a convenience sample of university students. Both factors limit generalisability: the system's effectiveness in a more complex commercial game with a more diverse player base remains an open question.

\textbf{Statistical conclusion validity.} With $N = 50$ and Bonferroni correction across three hypotheses ($\alpha_{\text{adj}} = 0.0167$), the study is not robustly powered for small effects. A true but small effect ($d < 0.4$) would likely fail to reach significance, leading to a Type II error. This limitation is acknowledged in the design and motivates the inclusion of qualitative data as a complementary source of evidence.

